\documentclass{beamer}

\usepackage{beamerthemeSingapore}
\usepackage[brazil]{babel}
\usepackage[T1]{fontenc}
\usepackage[utf8]{inputenc}
\usepackage{graphicx}
\usepackage{url}
\usepackage{xcolor}
\usepackage{booktabs}
\usepackage{ragged2e}
\usepackage{etoolbox}
\usepackage{minted}
\usepackage{array}

\newcommand{\tcb}[1]{\textcolor{blue}{#1}}
\newcommand{\tcr}[1]{\textcolor{red}{#1}}
\newcommand{\tcg}[1]{\textcolor{green!50!black}{#1}}

\setminted{
  fontsize=\footnotesize,
  breaklines=true,
  breakanywhere=true,
  frame=lines,
  framesep=2mm,
  baselinestretch=1.05,
  tabsize=2
}

\apptocmd{\frame}{}{\justifying}{}

\setbeamertemplate{navigation symbols}{}

% slides de prática e de resposta: fundo e título em cores próprias
\newcommand{\estilopratica}{%
  \setbeamercolor{frametitle}{fg=orange!85!black}%
  \setbeamercolor{background canvas}{bg=orange!8}}
\newcommand{\estiloresposta}{%
  \setbeamercolor{frametitle}{fg=green!45!black}%
  \setbeamercolor{background canvas}{bg=green!6}}

\title{07 - Fetch API e JSON}
\author{Prof. Alex Kutzke}
\date{Outubro de 2026}

\begin{document}

% ==========================================
\begin{frame}
	\begin{center}
		\large{Setor de Educação Profissional e Tecnológica}\\
		Tecnologia em Análise e Desenvolvimento de Sistemas\\
		\vspace{1cm}
		\small{DS122 - Desenvolvimento de Aplicações Web 1}
	\end{center}
	\titlepage
\end{frame}

% ==========================================
\begin{frame}{Objetivos da aula}
	Ao final desta aula você deve ser capaz de:
	\begin{enumerate}
		\item Escrever dados em JSON, e apontar o que muda em relação a um objeto
		      JavaScript;
		\item Converter entre texto JSON e valor com \texttt{JSON.parse} e
		      \texttt{JSON.stringify};
		\item Servir as páginas por um servidor local, e explicar por que o
		      \texttt{fetch} falha em \texttt{file://};
		\item Descrever o que \texttt{fetch} devolve, e usar \texttt{await} dentro
		      de uma função \texttt{async};
		\item Conferir \texttt{response.ok} e explicar por que um 404 não cai no
		      \texttt{catch};
		\item Carregar o catálogo de um arquivo JSON;
		\item Consultar uma API pública e tratar a resposta de erro dela.
	\end{enumerate}
\end{frame}

% ==========================================
\begin{frame}{Roteiro}
	\tableofcontents
\end{frame}

% ==========================================
\section{Requisição}
% ==========================================

\begin{frame}[fragile]{Uma requisição pelo console}
	\texttt{https://viacep.com.br} aberto, \texttt{F12}, aba \textbf{Network}
	aberta, depois aba \textbf{Console}:

	\begin{minted}{javascript}
const resposta = await fetch("https://viacep.com.br/ws/80060000/json/");
resposta.status        // 200
const dados = await resposta.json();
dados.logradouro       // "Rua XV de Novembro"
dados.localidade       // "Curitiba"
	\end{minted}

	\vspace{0.3cm}
	A página não recarregou.
\end{frame}

\begin{frame}{O que apareceu na aba Network}
	Uma linha nova, com:
	\begin{itemize}
		\item o endereço pedido;
		\item o método \texttt{GET};
		\item o status \texttt{200};
		\item na aba \textbf{Resposta}, um texto entre chaves, com pares de nome e
		      valor.
	\end{itemize}

	\vspace{0.3cm}
	O JavaScript fez uma requisição HTTP, como as da aula de HTTP, e transformou
	o corpo da resposta em objeto.

	\vspace{0.3cm}
	Nesta aula: de onde a página pode fazer requisições, o formato do corpo, por
	que o \texttt{await}, e o \texttt{fetch} em si.
\end{frame}

% ==========================================
\section{Servidor}
% ==========================================

\begin{frame}[fragile]{Página aberta do disco}
	Dois cliques no \texttt{catalogo.html}: o endereço começa com
	\texttt{file://}. Sem servidor e sem HTTP.

	\vspace{0.3cm}
	Para HTML, CSS e imagens, basta. Para o \texttt{fetch}, o navegador bloqueia:

	\begin{minted}[fontsize=\scriptsize]{text}
Access to fetch at 'file:///.../produtos.json' from origin 'null'
has been blocked by CORS policy: Cross origin requests are only
supported for protocol schemes: chrome, chrome-extension,
chrome-untrusted, data, http, https, isolated-app.
	\end{minted}

	\vspace{0.2cm}
	A partir desta aula, as páginas são abertas por um servidor local, em
	\texttt{http://localhost}.
\end{frame}

\begin{frame}[fragile]{Servidor local}
	Duas opções:

	\begin{itemize}
		\item VS Code: extensão \textbf{Live Server}, botão \emph{Go Live} no canto
		      inferior direito. Recarrega a página a cada alteração salva.
		\item Terminal, dentro da pasta do repositório:
	\end{itemize}

	\begin{minted}{bash}
python3 -m http.server 8000
	\end{minted}

	\quad e no navegador:

	\quad\texttt{http://localhost:8000/catalogo.html}

	\quad \texttt{Ctrl+C} encerra o servidor.

	\vspace{0.3cm}
	A aba \textbf{Network} passa a mostrar cada arquivo da página como uma
	requisição, com o status de cada um.
\end{frame}

{\estilopratica
\begin{frame}[fragile]{Prática: o ambiente}
	No seu \emph{fork} de \texttt{ds122-lab-02-assignment}:

	\begin{enumerate}
		\item copie suas páginas, \texttt{css/}, \texttt{imagens/},
		      \texttt{js/produtos.js} e \texttt{js/app.js} da tarefa de DOM, ou
		      mova o conteúdo de \texttt{base/} para a raiz;
		\item inicie o servidor local;
		\item abra o catálogo por \texttt{http://localhost};
		\item na aba \textbf{Network}, confira o status de cada arquivo.
	\end{enumerate}

	\vspace{0.3cm}
	Tempo: 5 minutos.
\end{frame}
}

{\estiloresposta
\begin{frame}[fragile]{Resposta: o ambiente}
	Os cartões aparecem, e a aba \textbf{Network} lista:

	\begin{minted}{text}
catalogo.html     200
estilo.css        200
produtos.js       200
app.js            200
produto.svg       200
	\end{minted}

	\vspace{0.2cm}
	Um arquivo com \texttt{404} é um caminho errado no HTML.
\end{frame}
}

% ==========================================
\section{JSON}
% ==========================================

\begin{frame}[fragile]{JSON}
	\emph{JavaScript Object Notation}: formato de texto para dados, lido e escrito
	por praticamente qualquer linguagem.

	\begin{minted}[fontsize=\scriptsize]{json}
[
  {
    "id": 1,
    "nome": "Café em grão",
    "preco": 32.5,
    "disponivel": true,
    "tags": ["orgânico", "torra média"]
  }
]
	\end{minted}

	Um único valor: objeto, array, string, número, \texttt{true},
	\texttt{false} ou \texttt{null}, um dentro do outro.
\end{frame}

\begin{frame}{O que muda em relação a um objeto JavaScript}
	\small
	\begin{tabular}{>{\raggedright\arraybackslash}p{0.45\textwidth}>{\raggedright\arraybackslash}p{0.45\textwidth}}
		\toprule
		No JavaScript & No JSON \\
		\midrule
		nome do campo com ou sem aspas & sempre entre aspas duplas \\ \midrule
		aspas simples, duplas ou crase & só aspas duplas \\ \midrule
		vírgula depois do último item & erro \\ \midrule
		comentários & não há \\ \midrule
		funções, \texttt{undefined} & não existem \\
		\bottomrule
	\end{tabular}

	\vspace{0.3cm}
	\texttt{"preco": "32.5"} é uma string. \texttt{"preco": 32.5} é um número, e
	só ele soma.
\end{frame}

\begin{frame}[fragile]{\texttt{JSON.parse} e \texttt{JSON.stringify}}
	O JSON chega pela rede como texto.

	\begin{minted}{javascript}
const texto = '{"nome": "Café em grão", "preco": 32.5}';

const produto = JSON.parse(texto);
produto.preco * 2;     // 65

JSON.stringify({ nome: "Mel", preco: 24.9 });
// '{"nome":"Mel","preco":24.9}'
	\end{minted}

	\vspace{0.2cm}
	Texto com erro de sintaxe lança \texttt{SyntaxError}. Teste:

	\begin{minted}{javascript}
JSON.parse("{'nome': 'Mel'}");
	\end{minted}

	\vspace{0.2cm}
	Na prática, o \texttt{response.json()} faz essa conversão por nós.
\end{frame}

{\estilopratica
\begin{frame}[fragile]{Prática: os produtos em JSON}
	\begin{enumerate}
		\item crie o \texttt{produtos.json} na raiz, com os produtos do
		      \texttt{js/produtos.js};
		\item acrescente a cada produto o campo \texttt{descricao};
		\item abra \texttt{http://localhost:8000/produtos.json} no Firefox e
		      confira que ele mostra os dados formatados;
		\item apague o \texttt{js/produtos.js} e a tag \texttt{<script>} dele.
	\end{enumerate}

	\vspace{0.3cm}
	Confira que a grade ficou vazia, com um erro no console.

	\vspace{0.3cm}
	Tempo: 10 minutos.
\end{frame}
}

{\estiloresposta
\begin{frame}[fragile]{Resposta: os produtos em JSON}
	\begin{minted}[fontsize=\scriptsize]{json}
[
  {
    "id": 1,
    "nome": "Café em grão",
    "descricao": "Grãos de torra média, de produção familiar.",
    "categoria": "Bebidas",
    "preco": 32.5,
    "desconto": 0.15,
    "imagem": "imagens/produto.svg"
  }
]
	\end{minted}

	\begin{minted}{text}
ReferenceError: produtos is not defined
	\end{minted}

	O \texttt{renderiza(produtos)} do fim do \texttt{app.js} ainda procura o
	array.
\end{frame}
}

% ==========================================
\section{Assíncrono}
% ==========================================

\begin{frame}{Operações assíncronas}
	Uma requisição leva tempo: milissegundos no servidor local, segundos numa
	conexão ruim.

	\vspace{0.3cm}
	Se o JavaScript parasse até a resposta chegar, a página congelaria: sem
	rolar, sem responder a cliques.

	\vspace{0.3cm}
	Por isso o \texttt{fetch} dispara a requisição e devolve, na hora, uma
	\emph{Promise} (promessa): um objeto que representa um resultado que ainda
	vai chegar.
\end{frame}

\begin{frame}[fragile]{Promise}
	\begin{minted}{javascript}
const p = fetch("https://viacep.com.br/ws/80060000/json/");
p   // Promise {<pending>}
	\end{minted}

	\vspace{0.2cm}
	Três estados:
	\begin{itemize}
		\item pendente (\emph{pending}): a resposta ainda não chegou;
		\item cumprida (\emph{fulfilled}): a resposta está dentro;
		\item rejeitada (\emph{rejected}): a requisição falhou, e o erro está
		      dentro.
	\end{itemize}
\end{frame}

\begin{frame}[fragile]{\texttt{await}}
	\texttt{await} antes de uma \emph{Promise} espera que ela se resolva e
	devolve o valor que está dentro:

	\begin{minted}{javascript}
const resposta = await fetch("https://viacep.com.br/ws/80060000/json/");
resposta   // Response {status: 200, ok: true, ...}
	\end{minted}

	\vspace{0.3cm}
	Fora do console, \texttt{await} só pode ser usado dentro de uma função
	marcada com \texttt{async}.
\end{frame}

% ==========================================
\section{Fetch}
% ==========================================

\begin{frame}{O objeto \texttt{Response}}
	\small
	\begin{tabular}{>{\raggedright\arraybackslash}p{0.35\textwidth}>{\raggedright\arraybackslash}p{0.55\textwidth}}
		\toprule
		Propriedade ou método & O que é \\
		\midrule
		\texttt{resposta.status} & código de status HTTP, como \texttt{200} ou
		\texttt{404} \\ \midrule
		\texttt{resposta.ok} & \texttt{true} com status entre 200 e 299 \\
		\midrule
		\texttt{resposta.json()} & lê o corpo e converte de JSON; devolve uma
		\emph{Promise} \\ \midrule
		\texttt{resposta.text()} & lê o corpo como texto; devolve uma
		\emph{Promise} \\
		\bottomrule
	\end{tabular}

	\vspace{0.3cm}
	O endereço pode ser completo, como o do ViaCEP, ou relativo à página, como
	\texttt{"produtos.json"}.
\end{frame}

\begin{frame}[fragile]{\texttt{response.ok} e o 404}
	O \texttt{fetch} só rejeita quando a requisição não acontece: sem rede,
	servidor fora do ar, domínio que não existe, bloqueio do navegador.

	\vspace{0.2cm}
	Um 404 é uma resposta HTTP completa:

	\begin{minted}{javascript}
const r = await fetch("nao-existe.json");
r.status   // 404
r.ok       // false
	\end{minted}

	\vspace{0.2cm}
	O console mostra a linha do 404 em vermelho, mas é um registro de rede do
	navegador: nenhum erro de JavaScript, e o código segue. Todo
	\texttt{fetch} confere \texttt{response.ok} antes de usar o corpo.
\end{frame}

\begin{frame}[fragile]{\texttt{response.json()}}
	Ler o corpo também é assíncrono: \texttt{json()} devolve outra
	\emph{Promise}.

	\begin{minted}{javascript}
const resposta = await fetch("produtos.json");
const produtos = await resposta.json();
produtos.length   // 6
	\end{minted}

	\vspace{0.3cm}
	Sem o segundo \texttt{await}, \texttt{produtos} é uma \emph{Promise}:

	\begin{minted}{text}
TypeError: produtos.forEach is not a function
	\end{minted}
\end{frame}

\begin{frame}[fragile]{Uma função \texttt{async} com \texttt{try} e \texttt{catch}}
	\begin{minted}[fontsize=\scriptsize]{javascript}
const carregaProdutos = async () => {
  try {
    const resposta = await fetch("produtos.json");
    if (!resposta.ok) {
      throw new Error(`Erro HTTP: ${resposta.status}`);
    }
    const produtos = await resposta.json();
    renderiza(produtos);
  } catch (erro) {
    console.error(erro);
    const aviso = document.createElement("p");
    aviso.classList.add("erro");
    aviso.textContent = "Não foi possível carregar os produtos.";
    grade.replaceChildren(aviso);
  }
};

carregaProdutos();
	\end{minted}
\end{frame}

\begin{frame}{Linha a linha}
	\begin{itemize}
		\item \texttt{async} marca a função como assíncrona e permite
		      \texttt{await} dentro dela;
		\item o \texttt{try} contém o caminho normal; qualquer erro desvia para o
		      \texttt{catch};
		\item \texttt{throw new Error(...)} lança um erro de propósito, e o 404
		      cai no mesmo \texttt{catch} que a falha de rede;
		\item o \texttt{catch} registra no console e avisa na página, porque o
		      visitante não abre o console;
		\item \texttt{renderiza} é a da aula de DOM, e não sabe de onde a lista
		      veio.
	\end{itemize}
\end{frame}

{\estilopratica
\begin{frame}[fragile]{Prática: o catálogo pelo \texttt{fetch}}
	\begin{enumerate}
		\item no fim do \texttt{app.js}, troque \texttt{renderiza(produtos);} por
		      uma \texttt{carregaProdutos} com \texttt{fetch}, \texttt{await} e
		      \texttt{response.ok}, ainda sem \texttt{try}/\texttt{catch};
		\item na \texttt{criaCartao}, acrescente um
		      \texttt{<p class="descricao">} logo depois da imagem;
		\item acrescente um sétimo produto ao \texttt{produtos.json}.
	\end{enumerate}

	\vspace{0.3cm}
	Confira que os cartões voltam com a descrição, que o sétimo aparece sem
	mexer no JavaScript, e que a aba \textbf{Network} mostra o
	\texttt{produtos.json} com \texttt{200}.

	\vspace{0.3cm}
	Tempo: 15 minutos.
\end{frame}
}

{\estiloresposta
\begin{frame}[fragile]{Resposta: o catálogo pelo \texttt{fetch}}
	\begin{minted}[fontsize=\scriptsize]{javascript}
const carregaProdutos = async () => {
  const resposta = await fetch("produtos.json");
  if (!resposta.ok) {
    throw new Error(`Erro HTTP: ${resposta.status}`);
  }
  const produtos = await resposta.json();
  renderiza(produtos);
};

carregaProdutos();
	\end{minted}

	\begin{minted}[fontsize=\scriptsize]{javascript}
// na criaCartao, depois da imagem
const descricao = document.createElement("p");
descricao.classList.add("descricao");
descricao.textContent = produto.descricao;
cartao.append(descricao);
	\end{minted}
\end{frame}
}

{\estilopratica
\begin{frame}[fragile]{Prática: a mensagem de erro}
	\begin{enumerate}
		\item envolva o corpo da \texttt{carregaProdutos} em
		      \texttt{try}/\texttt{catch};
		\item no \texttt{catch}, mostre
		      \texttt{Não foi possível carregar os produtos.} em um
		      \texttt{<p class="erro">} dentro da grade;
		\item troque \texttt{"produtos.json"} por \texttt{"produto.json"} no
		      \texttt{fetch} e recarregue.
	\end{enumerate}

	\vspace{0.3cm}
	Confira a mensagem na página, o erro no console e o \texttt{404} na aba
	\textbf{Network}. Desfaça a troca.

	\vspace{0.3cm}
	Tempo: 10 minutos.
\end{frame}
}

{\estiloresposta
\begin{frame}[fragile]{Resposta: a mensagem de erro}
	O código do slide \emph{Uma função \texttt{async} com \texttt{try} e
	\texttt{catch}}.

	\vspace{0.3cm}
	Com \texttt{"produto.json"}:

	\begin{minted}{text}
Network:  produto.json   404
Console:  Error: Erro HTTP: 404
Página:   Não foi possível carregar os produtos.
	\end{minted}

	\vspace{0.2cm}
	Sem o \texttt{if (!resposta.ok)}, o \texttt{json()} tenta ler a página de
	erro do servidor como JSON, e o \texttt{catch} recebe um
	\texttt{SyntaxError} que não diz que o arquivo não foi encontrado.
\end{frame}
}

\begin{frame}[fragile]{A ordem de execução}
	\begin{minted}[fontsize=\scriptsize]{javascript}
let produtos = [];

const carrega = async () => {
  const resposta = await fetch("produtos.json");
  produtos = await resposta.json();
  console.log("2. dados chegaram:", produtos.length);
};

carrega();
console.log("1. logo depois da chamada:", produtos.length);
	\end{minted}

	\begin{minted}{text}
1. logo depois da chamada: 0
2. dados chegaram: 6
	\end{minted}

	No primeiro \texttt{await}, a função devolve o controle. O que usa os dados
	vai \textbf{dentro} dela, depois do \texttt{await}.
\end{frame}

% ==========================================
\section{API}
% ==========================================

\begin{frame}[fragile]{Uma API pública: ViaCEP}
	API, na web: um endereço que responde dados em vez de páginas, em geral em
	JSON.

	\begin{minted}{text}
https://viacep.com.br/ws/80060000/json/
	\end{minted}

	\small
	\begin{tabular}{>{\raggedright\arraybackslash}p{0.38\textwidth}>{\raggedright\arraybackslash}p{0.12\textwidth}>{\raggedright\arraybackslash}p{0.35\textwidth}}
		\toprule
		Pedido & Status & Corpo \\
		\midrule
		\texttt{123}, formato inválido & 400 & página de erro em HTML \\ \midrule
		\texttt{99999999}, não existe & 200 &
		\texttt{\{ "erro": "true" \}} \\
		\bottomrule
	\end{tabular}

	\vspace{0.3cm}
	No segundo caso, \texttt{response.ok} é \texttt{true}. O erro está no
	conteúdo, e a documentação da API diz como.
\end{frame}

\begin{frame}[fragile]{Formato do CEP e CORS}
	O formato pode ser conferido antes da requisição:

	\begin{minted}{javascript}
const cep = campoCep.value.replace("-", "");
if (cep.length !== 8 || Number.isNaN(Number(cep))) {
  // mostra a mensagem e não faz o fetch
}
	\end{minted}

	\vspace{0.2cm}
	O navegador só entrega a resposta de outro site se o servidor permitir, pelo
	cabeçalho \texttt{Access-Control-Allow-Origin}. É o CORS
	(\emph{Cross-Origin Resource Sharing}), o mesmo mecanismo que bloqueia
	\texttt{file://}. O ViaCEP permite qualquer site.
\end{frame}

{\estilopratica
\begin{frame}[fragile]{Prática: o endereço pelo CEP}
	\begin{enumerate}
		\item no \texttt{contato.html}, acrescente os campos \texttt{cep},
		      \texttt{rua}, \texttt{bairro} e \texttt{cidade}, e um
		      \texttt{<p id="erro-cep" class="erro">};
		\item crie o \texttt{js/contato.js}, ligado com \texttt{defer};
		\item escreva a \texttt{buscaCep}, chamada no \texttt{blur} do
		      \texttt{cep}, que confere o formato, faz o \texttt{fetch} e preenche
		      os três campos.
	\end{enumerate}

	\vspace{0.3cm}
	Confira com \texttt{80060000}, \texttt{123} e \texttt{99999999}.

	\vspace{0.3cm}
	Tempo: 10 minutos.
\end{frame}
}

{\estiloresposta
\begin{frame}[fragile]{Resposta: o endereço pelo CEP}
	\begin{minted}[fontsize=\scriptsize]{javascript}
const campoCep = document.querySelector("#cep");
const erroCep = document.querySelector("#erro-cep");
const buscaCep = async () => {
  erroCep.textContent = "";
  const cep = campoCep.value.replace("-", "");
  if (cep.length !== 8 || Number.isNaN(Number(cep))) {
    erroCep.textContent = "Digite um CEP com oito números.";
    return;
  }
  const resposta = await fetch(`https://viacep.com.br/ws/${cep}/json/`);
  const dados = await resposta.json();
  document.querySelector("#rua").value = dados.logradouro;
  document.querySelector("#bairro").value = dados.bairro;
  document.querySelector("#cidade").value = dados.localidade;
};
campoCep.addEventListener("blur", buscaCep);
	\end{minted}

	Com \texttt{99999999}, os campos ficam com \texttt{undefined}. O tratamento
	do \texttt{erro} e o \texttt{try}/\texttt{catch} ficam para a tarefa.
\end{frame}
}

\begin{frame}{Quando dá errado}
	\small
	\begin{tabular}{>{\raggedright\arraybackslash}p{0.45\textwidth}>{\raggedright\arraybackslash}p{0.45\textwidth}}
		\toprule
		Sintoma & Causa provável \\
		\midrule
		\texttt{blocked by CORS policy}, com \texttt{file://} & página sem
		servidor local \\ \midrule
		\texttt{SyntaxError} ao ler o JSON & vírgula sobrando, aspas simples,
		comentário \\ \midrule
		\texttt{.forEach is not a function} & \texttt{await} esquecido no
		\texttt{json()} \\ \midrule
		\texttt{await is only valid in async functions} & \texttt{await} fora
		de função \texttt{async} \\ \midrule
		\texttt{Unexpected token \textquotesingle<\textquotesingle}, com 404 & \texttt{json()} lendo a página de
		erro, sem conferir \texttt{response.ok} \\ \midrule
		total ou contagem zerados & dados usados fora da função \\ \midrule
		soma dá \texttt{"032.518"} & números entre aspas no JSON \\
		\bottomrule
	\end{tabular}

	\vspace{0.2cm}
	Os dois últimos não aparecem no console. A aba \textbf{Network} mostra o
	status e o corpo que chegou.
\end{frame}

\begin{frame}{Tarefa}
	Os exercícios valem nota e compõem o item \emph{Exercícios em sala}.

	\vspace{0.3cm}
	Enunciado no \texttt{README.md} de \texttt{ds122-lab-02-assignment}; prazo
	na UFPR Virtual. As paradas desta aula cobrem os exercícios 1 a 4; o
	exercício 5 é de revisão para a Prova 2.

	\vspace{0.3cm}
	Entrega pelo \emph{fork} no GitLab, individual ou em dupla, sem uso de IA
	generativa para produzir o código.

	\vspace{0.5cm}
	O catálogo pelo \texttt{produtos.json} é o requisito 4 da Entrega 2, com
	prazo no domingo, 11/10. No trabalho, o cartão mantém o que a Entrega 1
	pedia: imagem, nome, descrição curta e preço.
\end{frame}

\end{document}
